% Nouveau EGA V (cote 25): archive cote 25, site page 7 (PDF page 8).
% typescript. Draft compared: none.
% AI-transcribed by Claude Opus 5.5 (claude-opus-5-5, adaptive thinking, effort high), Message Batches API, batch ega25b (msgbatch_01A7n5YuXubUPNK1zsepi8RG), 2026-11-13; original AI transcription, subsequently checked against both supplied scan parts; see PAGES_CHECKED.tsv and CHECK_DECISIONS.tsv.
\scanpage{25/7}
\mspage{\hand{4}}

par le morphisme composé $H_{\hand{$\xi$}} \hand{\longrightarrow} P_{S'} \hand{\longrightarrow} P$. Il est alors commode
\xout{\illegible{4}} encore de considérer le cas où \hand{$\xi$} est
``universel'' i.e. où $S' = \overset{\hand{\vee}}{P}$ , et \hand{$\xi$} est la section diagonale, donc $H \ = \underline{H}$,
dans ce cas on notera (sauf meilleures notations suggérées par Dieudonné)
$\underline{Y} = Y_{\hand{$\xi$}}$ . Dans le cas général d'un \hand{$\xi$}: $S' \hand{\to} \overset{\hand{\vee}}{P}$ , on aura donc aussi
\[ Y_{\hand{$\xi$}} = \underline{Y} \times_{\overset{\hand{\vee}}{P}} S' \ . \]
Enfin, si F est Module sur X, nous \unsure{déssignerons|désignerons} par $G_{\hand{$\xi$}}$ son image invers
% NOTE: the right edge of the page is cut off here; the end of "inverse" is not visible
sur $Y_{\hand{$\xi$}}$ , par $\underline{G}$ son image inverse sur $\unsure{\underline{M}|\underline{H}}$ , de sorte qu'on aura aussi
\[ G_{\hand{$\xi$}} = \underline{G} \otimes_{\underline{O}_{\overset{\hand{\vee}}{P}}} \underline{O}_{S'} \ . \]
Résumons dans un petit diagramme l'essent\unsure{i}el des constructions \inserted{et notations} envisag
% NOTE: the right edge of the page is cut off at the end of this line

gées:
\diagram{
\begin{tikzcd}[arrows={blue}]
F & & \underline{G} & & G_{\hand{$\xi$}} \\
X \arrow[d] & X\times_S\overset{\hand{\vee}}{P} \arrow[l] \arrow[d] & \underline{Y} \arrow[l] \arrow[d] & & Y_{\hand{$\xi$}} \arrow[ll] \arrow[d] \\
P \arrow[d] & P\times_S\overset{\hand{\vee}}{P} \arrow[l] \arrow[d] & \underline{H} \arrow[l] \arrow[dl] & & H_{\hand{$\xi$}} \arrow[ll] \arrow[dl] \\
S & \overset{\hand{\vee}}{P} \arrow[l] & & S' \arrow[ll] &
\end{tikzcd}
}
\margin{\illegible{1}}

(Les carrés et losanges figurant dans ce diagramme sont car\over{\illegible{1}}{té}siens).

Dans le numéro suivant, nous étudions systématiquement le cas
s\unsure{u}ivant: S' est le spectre d'un corps\inserted{K}, et son image dahs P est \uline{génériqu}
% NOTE: the right edge of the page is cut off at the end of this line
dans la fibre correspondante $\overset{\hand{\vee}}{P}_s$ . Quitte à faire le changement de base
Spec k(s) $\hand{\longrightarrow}$ S , on est ramené pour ce\xout{\illegible{1}} au cas où S est le spectre
d'un corps k, ce que nous supposerons au N\textsuperscript{o} suivant. D'ailleurs, la plu-%
part des propriétés examinées pour X et $Y_{\hand{$\xi$}}$ sont de ``nature \unsure{géométriqque|géométrique}''
donc invariantes par changement de corps de base; cela permet alors
indifféremment de se ramener au cas où K est algébriquement clos, ou au
cas où K = k(\hand{$\eta$}) , \hand{$\eta$} étant le point générique de $\overset{\hand{\vee}}{P}$ , et \hand{$\xi$} : Spec(K)$\hand{\to}\overset{\hand{\vee}}{P}$
étant bien entendu le morphisme canonique. On notera de même que pour le
% NOTE: the right edge of the page is cut off at the end of this line
questions de nature géométrique concernant X, $Y_{\hand{$\xi$}}$ , on peut (quitte à fa
% NOTE: the right edge of the page is cut off at the end of this line
re un changement de base sur k) se borner au cas k alg clos.

\mspage{\hand{7}}
